%=====================================================================
\chapter{Clustering}\label{ch:clustering}
%=====================================================================
\locator{4}{Part IV --- Learning Without Full Labels}

\begin{outcomes}
By the end of this chapter you will be able to:
\begin{tight}
  \item \textbf{Explain} what clustering is for, and why its answer depends on
        choices the analyst makes.
  \item \textbf{Trace} $k$-means by hand through assignment and update steps to
        convergence.
  \item \textbf{Build} a dendrogram by hierarchical agglomerative clustering under
        single, complete and average linkage.
  \item \textbf{Choose} the number of clusters with the elbow method and the
        silhouette coefficient.
  \item \textbf{Evaluate} a clustering internally and, when labels exist,
        externally.
\end{tight}
\end{outcomes}

\begin{prereq}
\chref{math} (distance), \chref{workflow} (scaling) and \chref{knn} (distance
measures and the curse of dimensionality, which clustering inherits in full).
\end{prereq}

\begin{keyterms}
clustering $\bullet$ partitional vs.\ hierarchical $\bullet$ centroid $\bullet$
$k$-means (Lloyd's algorithm) $\bullet$ within-cluster sum of squares (WCSS)
$\bullet$ local minimum $\bullet$ $k$-means++ $\bullet$ agglomerative clustering
$\bullet$ dendrogram $\bullet$ single, complete, average and Ward linkage
$\bullet$ chaining $\bullet$ elbow method $\bullet$ silhouette coefficient
$\bullet$ adjusted Rand index
\end{keyterms}

\section{Finding Groups Without Labels}

\term{Clustering} divides examples into groups --- \term{clusters} --- so that
examples in the same group are similar and examples in different groups are not.
There are no labels to learn from and no right answer to check against. What
counts as ``similar'' is a choice (of features, their scaling and a distance), and
so is the number of groups; change either and the clusters change.

That makes clustering exploratory: a way of proposing structure for a person to
examine. Customer segments, kinds of network traffic, families of malware, modes
in which a machine operates --- each is a clustering that someone then named and
decided was useful. Two families of algorithm cover the HEC outline:
\term{partitional} methods such as $k$-means divide the data into a fixed number of
groups; \term{hierarchical} methods build a whole tree of nested groupings.

\section{$k$-Means}

\begin{definitionbox}[$k$-Means (Lloyd's Algorithm)]
Choose $k$ and place $k$ initial \term{centroids}. Repeat until nothing changes:
\begin{tightnum}
  \item \textbf{Assign} each example to its nearest centroid.
  \item \textbf{Update} each centroid to the mean of the examples assigned to it.
\end{tightnum}
Each step can only lower the \term{within-cluster sum of squares}
\[
\text{WCSS} = \sum_{j=1}^{k}\ \sum_{\mathbf{x} \in C_j} \|\mathbf{x} - \boldsymbol{\mu}_j\|^2
\]
so the algorithm always converges --- but to a \term{local} minimum that depends
on where it started.
\end{definitionbox}

\begin{worked}[$k$-Means by Hand]
Seven points: A1 $(1, 1)$, A2 $(1.5, 2)$, A3 $(3, 4)$, A4 $(5, 7)$, A5 $(3.5, 5)$,
A6 $(4.5, 5)$, A7 $(3.5, 4.5)$. $k = 2$, with initial centroids A2 and A4.

\smallskip
\textbf{Iteration 1 --- assign.} Distances to $\mu_1 = (1.5, 2)$ and $\mu_2 = (5, 7)$:
A1 1.12 / 7.21; A2 0 / 6.10; A3 2.50 / 3.61; A4 6.10 / 0; A5 3.61 / 2.50;
A6 4.24 / 2.06; A7 3.20 / 2.92. So $C_1 = \{$A1, A2, A3$\}$, $C_2 = \{$A4, A5, A6, A7$\}$.

\textbf{Update.} $\mu_1 = \big(\tfrac{1+1.5+3}{3}, \tfrac{1+2+4}{3}\big) = (1.83, 2.33)$;
$\mu_2 = \big(\tfrac{5+3.5+4.5+3.5}{4}, \tfrac{7+5+5+4.5}{4}\big) = (4.13, 5.38)$.
WCSS $= 12.21$.

\smallskip
\textbf{Iteration 2 --- assign.} A3 is now 2.03 from $\mu_1$ and \textbf{1.78} from
$\mu_2$: it moves to $C_2$. Every other point stays.

\textbf{Update.} $\mu_1 = (1.25, 1.5)$, $\mu_2 = (3.9, 5.1)$. WCSS $= 8.53$.

\smallskip
\textbf{Iteration 3.} No point changes cluster, so the centroids stop moving:
\textbf{converged}, with $C_1 = \{$A1, A2$\}$ and $C_2 = \{$A3, \dots, A7$\}$.
\end{worked}

\dsfig{%
\begin{tikzpicture}[font=\scriptsize, scale=0.62,
  c1/.style={circle, fill=secondary, inner sep=1.7pt},
  c2/.style={circle, fill=accent, inner sep=1.7pt},
  cen/.style={cross out, draw=#1, line width=1.6pt, minimum size=7pt, inner sep=0pt}]
\foreach \ox/\ttl in {0/{Start: centroids at A2, A4}, 7.2/{After iteration 1}, 14.4/{Converged (iteration 2)}}{
  \begin{scope}[shift={(\ox,0)}]
  \draw[gray!50, -{Stealth[length=1.5mm]}] (0,0) -- (6.2,0);
  \draw[gray!50, -{Stealth[length=1.5mm]}] (0,0) -- (0,7.8);
  \node[font=\scriptsize\bfseries, text=primary] at (3.1,-0.8) {\ttl};
  \end{scope}}
% panel 1: initial assignment
\foreach \x/\y/\s in {1/1/c1, 1.5/2/c1, 3/4/c1, 5/7/c2, 3.5/5/c2, 4.5/5/c2, 3.5/4.5/c2}{\node[\s] at (\x,\y) {};}
\node[cen=primary] at (1.5,2) {}; \node[cen=primary] at (5,7) {};
% panel 2
\begin{scope}[shift={(7.2,0)}]
\foreach \x/\y/\s in {1/1/c1, 1.5/2/c1, 3/4/c2, 5/7/c2, 3.5/5/c2, 4.5/5/c2, 3.5/4.5/c2}{\node[\s] at (\x,\y) {};}
\node[cen=primary] at (1.833,2.333) {}; \node[cen=primary] at (4.125,5.375) {};
\draw[goldhl, line width=1pt] (3,4) circle (0.32);
\node[font=\tiny, goldhl!45!black, anchor=west] at (3.35,3.55) {A3 switches};
\end{scope}
% panel 3
\begin{scope}[shift={(14.4,0)}]
\foreach \x/\y/\s in {1/1/c1, 1.5/2/c1, 3/4/c2, 5/7/c2, 3.5/5/c2, 4.5/5/c2, 3.5/4.5/c2}{\node[\s] at (\x,\y) {};}
\node[cen=primary] at (1.25,1.5) {}; \node[cen=primary] at (3.9,5.1) {};
\draw[secondary!60, dashed] (1.25,1.5) circle (0.9);
\draw[accent!60, dashed] (3.9,5.1) circle (2.1);
\end{scope}
\node[anchor=west, font=\tiny, text=gray!55!black] at (0,-1.6)
     {colour = current cluster; $\times$ = centroid. Panel 2 shows the assignment made with the
      updated centroids, in which A3 changes cluster.};
\end{tikzpicture}}%
{Three snapshots of $k$-means on the worked example. Assignment and update alternate until
no point changes cluster.}{kmeans-steps}

\begin{alertbox}[Where $k$-Means Goes Wrong]
\begin{tight}
  \item \textbf{Bad starts.} Different initial centroids can converge to different,
        worse clusterings. Run it many times and keep the lowest WCSS
        (\texttt{n\_init}), and start with \term{$k$-means++}, which spreads the
        initial centroids out.
  \item \textbf{Its shape bias.} Every cluster is the set of points nearest a
        centroid, so $k$-means assumes round, similar-sized clusters. Long,
        curved or very unequal clusters are cut badly.
  \item \textbf{Scale.} It is built on Euclidean distance; unscaled features
        decide everything (\chref{knn}).
  \item \textbf{Outliers.} A mean is pulled by extreme values; one wild point can
        capture a centroid of its own.
\end{tight}
\end{alertbox}

\section{Hierarchical Agglomerative Clustering}

\begin{definitionbox}[Agglomerative Clustering and Linkage]
Start with every example in a cluster of its own. Repeatedly merge the two
\emph{closest} clusters until one remains. The record of merges, drawn as a tree
whose heights are the merge distances, is a \term{dendrogram}. The distance between
two clusters is defined by the \term{linkage}:
\begin{tight}
  \item \term{single}: the distance between their \emph{closest} members;
  \item \term{complete}: the distance between their \emph{farthest} members;
  \item \term{average}: the average distance over all pairs of members;
  \item \term{Ward}: the increase in WCSS that merging them would cause --- the
        $k$-means criterion, applied hierarchically.
\end{tight}
\end{definitionbox}

\begin{worked}[Single and Complete Linkage on Five Numbers]
Points on a line: 1, 2, 4, 9, 11.

\smallskip
\textbf{Single linkage.} Closest pair: \{1, 2\} at distance 1. Next, \{1, 2\} to 4 is
$|2 - 4| = 2$ and 9 to 11 is also 2: merge \{1, 2, 4\} and \{9, 11\}, both at height
2. Last, the two groups are $|4 - 9| = 5$ apart. \textbf{Heights: 1, 2, 2, 5.}

\textbf{Complete linkage.} \{1, 2\} at 1. Now \{1, 2\} to 4 is $|1 - 4| = 3$ (the
farthest pair), while 9 to 11 is 2, so \{9, 11\} merges first, at 2; then
\{1, 2, 4\} at 3. Last, the farthest pair across the groups is $|1 - 11| = 10$.
\textbf{Heights: 1, 2, 3, 10.}

\smallskip
Both find the same two groups, but complete linkage separates them by a much taller
final merge. Average linkage gives a final height of $(8 + 10 + 7 + 9 + 5 + 7)/6 =
7.67$. \textbf{Cutting} a dendrogram at any height yields a flat clustering: cut
complete linkage anywhere between 3 and 10 and you get \{1, 2, 4\} and \{9, 11\}.
\end{worked}

\dsfig{%
\begin{tikzpicture}[font=\scriptsize, yscale=0.32, xscale=0.9]
\foreach \ox/\ttl/\h/\hh/\top in {0/{Single linkage}/2/2/5, 7.2/{Complete linkage}/3/2/10}{
  \begin{scope}[shift={(\ox,0)}]
  \draw[gray!55, -{Stealth[length=1.5mm]}] (-0.5,0) -- (-0.5,11) node[above, font=\tiny] {height};
  \foreach \t in {0,2,4,6,8,10}{\draw[gray!40] (-0.6,\t) -- (-0.4,\t) node[left, font=\tiny] at (-0.6,\t) {\t};}
  \foreach \i/\v in {0/1,1/2,2/4,3/9,4/11}{\node[font=\scriptsize] at (\i,-0.9) {\v};}
  % {1,2} at 1
  \draw[primary, line width=1pt] (0,0) -- (0,1) -- (1,1) -- (1,0);
  % {1,2}+4 at \h
  \draw[primary, line width=1pt] (0.5,1) -- (0.5,\h) -- (2,\h) -- (2,0);
  % {9,11} at \hh
  \draw[accent, line width=1pt] (3,0) -- (3,\hh) -- (4,\hh) -- (4,0);
  % final
  \draw[gray!60!black, line width=1pt] (1.25,\h) -- (1.25,\top) -- (3.5,\top) -- (3.5,\hh);
  \node[font=\scriptsize\bfseries, text=primary] at (1.75,12.3) {\ttl};
  \end{scope}}
\draw[goldhl, dashed, line width=0.9pt] (6.4,6) -- (11.6,6);
\node[font=\tiny, goldhl!45!black, anchor=west] at (11.6,6) {cut: 2 clusters};
\end{tikzpicture}}%
{Dendrograms for the points 1, 2, 4, 9, 11. The merges happen in almost the same order, but
at different heights: complete linkage keeps groups compact and makes the gap between them
obvious.}{dendrogram}

\begin{conceptbox}[Chaining, and Which Linkage to Use]
Single linkage merges two clusters as soon as \emph{any} pair of their members is close,
so a thin line of points can join two genuinely separate groups into one long
``chain''. That lets it find elongated or curved clusters, and also makes it fragile
with noise. Complete linkage does the opposite: it favours compact, roughly equal
clusters and is pulled by outliers. Average linkage sits between them. Ward's method
behaves like $k$-means and is a good default. Unlike $k$-means, hierarchical
clustering needs no $k$ in advance and gives the same answer every run --- but it
must compute all $n^2$ pairwise distances, which limits it to tens of thousands of
examples.
\end{conceptbox}

\section{How Many Clusters?}

\textbf{The elbow method.} Plot WCSS against $k$. WCSS always falls as $k$ rises ---
with $k = n$ it is zero --- so look for the \emph{elbow}, the $k$ after which extra
clusters buy little.

\textbf{The silhouette.} For each example, let $a$ be its mean distance to the other
members of its own cluster and $b$ its mean distance to the members of the nearest
other cluster. Its \term{silhouette} is
\[
s = \frac{b - a}{\max(a, b)} \qquad (-1 \le s \le 1)
\]
Near 1: well inside its cluster. Near 0: on a boundary. Negative: probably in the
wrong cluster. The mean silhouette over all examples scores a whole clustering; pick
the $k$ that maximises it. A point with $a = 2$ and $b = 6$ has $s = (6 - 2)/6 = 0.67$.

\dsfig{%
\begin{tikzpicture}
\begin{groupplot}[
  group style={group size=2 by 1, horizontal sep=1.5cm},
  width=6.8cm, height=4.8cm, axis lines=left,
  tick label style={font=\tiny}, label style={font=\tiny},
  title style={font=\scriptsize\bfseries}, xlabel={number of clusters $k$},
  xmin=0.5, xmax=8.5, xtick={1,...,8}]
\nextgroupplot[title={Elbow: WCSS}, ylabel={WCSS (thousands)}, ymin=0, ymax=40]
\addplot[secondary, line width=1.2pt, mark=*, mark size=1.5pt] coordinates
  {(1,37.23) (2,9.25) (3,4.37) (4,1.21) (5,1.09) (6,0.98) (7,0.87) (8,0.77)};
\draw[accent, line width=0.9pt] (axis cs:4,1.21) circle (4pt);
\node[font=\tiny, accent, anchor=south west] at (axis cs:4.2,4) {elbow at $k = 4$};
\nextgroupplot[title={Mean silhouette}, ylabel={silhouette}, ymin=0.25, ymax=0.8]
\addplot[greenhl!70!black, line width=1.2pt, mark=*, mark size=1.5pt] coordinates
  {(2,0.636) (3,0.608) (4,0.728) (5,0.624) (6,0.521) (7,0.432) (8,0.325)};
\draw[accent, line width=0.9pt] (axis cs:4,0.728) circle (4pt);
\node[font=\tiny, accent, anchor=west] at (axis cs:4.3,0.745) {maximum at $k = 4$};
\end{groupplot}
\end{tikzpicture}}%
{Choosing $k$ for 600 points generated from four clusters. The WCSS curve bends sharply at
$k = 4$ and the silhouette peaks there: both agree with the truth.}{choose-k}

\begin{alertbox}[Judging a Clustering]
\textbf{Internal} measures (WCSS, silhouette) use only the data, and reward whatever
shape the measure assumes. \textbf{External} measures such as the \term{adjusted Rand
index} compare the clusters with known labels, when some exist: 1 for perfect
agreement, about 0 for chance. The most important test is neither: can a domain
expert name the clusters, and do the names lead to a different decision for each?
A clustering nobody can use is not a result.
\end{alertbox}

\section{Beyond the HEC Outline}

Two methods complete the picture. \textbf{DBSCAN} grows clusters from dense regions
and labels isolated points as noise, so it finds clusters of any shape without being
told $k$ --- it belongs to the Data Mining course, with outlier detection.
\textbf{Gaussian mixture models} replace $k$-means's hard assignments by
probabilities; they are fitted with the EM algorithm of \chref{semisup}, of which
$k$-means turns out to be a special case.

\begin{pitfall}
\begin{tight}
  \item \textbf{Clustering unscaled features.} The feature with the largest units
        decides the clusters.
  \item \textbf{Running $k$-means once.} Use many starts and $k$-means++.
  \item \textbf{Choosing $k$ by WCSS alone.} It always falls; look for the elbow and
        check the silhouette.
  \item \textbf{Treating clusters as discovered facts.} They reflect the features,
        scaling, distance and $k$ chosen; a different choice gives different
        clusters.
  \item \textbf{Forcing $k$-means onto non-round clusters.} Use single linkage,
        DBSCAN or a mixture model instead.
\end{tight}
\end{pitfall}

%---------------------------------------------------------------------
\begin{fourlenses}{Clustering}
\lensLA{Cluster students by how they study --- when they log in, how they use
videos and quizzes --- and advisers can tailor support to ``last-minute crammers''
or ``steady early starters''. The names come from people looking at the clusters,
and the clusters change if you add or remove a behavioural feature.}

\lensPM{Clustering past projects on size, duration, team and technology yields
project archetypes, each with its own typical overrun. A new project is estimated
from its archetype --- nearest-centroid classification built from an unsupervised
start.}

\lensIOT{Clustering a machine's sensor readings recovers its operating modes ---
idle, normal load, heavy load --- without anyone labelling them. A reading far from
every centroid is a candidate anomaly; one that drifts from its mode's centroid over
weeks is wear.}

\lensSEC{Clustering malware by behaviour groups new samples into families, so
analysts examine one sample per cluster instead of thousands. Network traffic
clusters separate normal services; a small, tight new cluster appearing overnight
deserves attention.}
\end{fourlenses}

%---------------------------------------------------------------------
\begin{lab}[k-Means, a Dendrogram and the Choice of k]
\begin{lstlisting}[language=Python]
import warnings
import numpy as np
from scipy.cluster.hierarchy import linkage, fcluster
from sklearn.datasets import make_blobs
from sklearn.cluster import KMeans, AgglomerativeClustering
from sklearn.metrics import silhouette_score, adjusted_rand_score
from sklearn.preprocessing import StandardScaler

warnings.filterwarnings("ignore", message="KMeans is known to have a memory")

# --- 1. the worked example: k-means from A2 and A4 -----------------------
P = np.array([[1, 1], [1.5, 2], [3, 4], [5, 7], [3.5, 5], [4.5, 5],
              [3.5, 4.5]])
km = KMeans(n_clusters=2, init=P[[1, 3]], n_init=1).fit(P)
print("centroids:", km.cluster_centers_.round(3).tolist(),
      " WCSS:", round(km.inertia_, 3), " iterations:", km.n_iter_)

# --- 2. choosing k: the elbow and the silhouette -------------------------
X, truth = make_blobs(n_samples=600, centers=4, cluster_std=1.0,
                      random_state=3)
for k in range(2, 8):
    m = KMeans(n_clusters=k, n_init=10, random_state=0).fit(X)
    print(f"k={k}: WCSS {m.inertia_:8.1f}   "
          f"silhouette {silhouette_score(X, m.labels_):.3f}")

# --- 3. hierarchical clustering: linkage heights on five numbers ---------
x = np.array([[1.0], [2.0], [4.0], [9.0], [11.0]])
for method in ("single", "complete", "average"):
    Z = linkage(x, method)
    print(f"{method:8s} merge heights: {Z[:, 2].round(2).tolist()}")
print("complete linkage cut into 2 groups:",
      fcluster(linkage(x, "complete"), 2, criterion="maxclust").tolist())

# --- 4. against the truth, and why scale matters -------------------------
km4 = KMeans(4, n_init=10, random_state=0).fit_predict(X)
hac = AgglomerativeClustering(n_clusters=4, linkage="ward").fit(X)
print("adjusted Rand, k-means:", round(adjusted_rand_score(truth, km4), 3),
      " ward HAC:", round(adjusted_rand_score(truth, hac.labels_), 3))
Xs = X.copy()
Xs[:, 1] *= 100                                   # one feature in other units
raw = KMeans(4, n_init=10, random_state=0).fit_predict(Xs)
fix = KMeans(4, n_init=10, random_state=0).fit_predict(
    StandardScaler().fit_transform(Xs))
print("rescaled feature, adjusted Rand: raw",
      round(adjusted_rand_score(truth, raw), 3), " standardised",
      round(adjusted_rand_score(truth, fix), 3))
\end{lstlisting}
The warning filter hides a harmless message scikit-learn prints on some Windows
installations. Change \texttt{cluster\_std} to 2.5 in part 2: do the elbow and the
silhouette still agree?
\end{lab}

%---------------------------------------------------------------------
\begin{checkpoint}
\begin{tightnum}
  \item Two centroids are at 2 and 8 on a line. Assign the points 1, 4, 6, 7, 9 and
        compute the new centroids.
  \item Why does $k$-means always converge, and why might it converge to a poor
        answer?
  \item A point has $a = 3$ and $b = 2$. Compute its silhouette and say what it
        means.
\end{tightnum}
\end{checkpoint}

%---------------------------------------------------------------------
\begin{chaptersummary}
\begin{tight}
  \item Clustering proposes groups without labels; the result depends on features,
        scaling, distance and the number of clusters.
  \item $k$-means alternates assignment to the nearest centroid with moving each
        centroid to its cluster's mean, lowering WCSS until it converges to a local
        minimum; use many starts and $k$-means++.
  \item Agglomerative clustering merges the closest clusters step by step into a
        dendrogram; the linkage (single, complete, average, Ward) defines
        ``closest''. Single linkage chains; complete linkage keeps groups compact.
  \item Choose $k$ with the elbow of the WCSS curve and the maximum mean
        silhouette; judge the result by whether it can be named and used.
\end{tight}
\end{chaptersummary}

\begin{reviewq}
  \item \textbf{(MCQ)} The $k$-means update step moves each centroid to:
        (a)~the median of its cluster (b)~the mean of its cluster (c)~the farthest
        point (d)~a random point.
  \item \textbf{(MCQ)} Which linkage is most prone to chaining? (a)~single
        (b)~complete (c)~average (d)~Ward.
  \item \textbf{(MCQ)} A silhouette close to $-1$ suggests the point is:
        (a)~well clustered (b)~on a boundary (c)~probably in the wrong cluster
        (d)~an outlier with no cluster.
  \item \textbf{(Short)} Why can the WCSS alone not be used to choose $k$?
  \item \textbf{(Short)} Give two situations where hierarchical clustering is
        preferable to $k$-means.
  \item \textbf{(Short)} Explain the role of $k$-means++.
  \item \textbf{(Applied)} Run $k$-means by hand on the worked example with initial
        centroids A1 and A7. Does it reach the same final clusters?
  \item \textbf{(Applied)} Build single- and complete-linkage dendrograms for the
        points 0, 3, 4, 10, 12, 20 and state the two-cluster cut of each.
\end{reviewq}
